%=====================================================================
\chapter{Keys and Integrity Constraints}\label{ch:constraints}
%=====================================================================
\locator{2}{Part II --- Designing It}

\begin{outcomes}
By the end of this chapter you will be able to:
\begin{tight}
  \item \textbf{State} the three model-based integrity constraints --- domain,
        entity and referential --- and \textbf{give} an example of each being
        violated.
  \item \textbf{Define} a foreign key and \textbf{explain} what referential
        integrity forbids.
  \item \textbf{Choose} an appropriate referential action (\sqlkw{CASCADE},
        \sqlkw{RESTRICT}, \sqlkw{SET NULL}) and \textbf{justify} the choice.
  \item \textbf{Demonstrate} that a declared constraint and an enforced
        constraint are not the same thing.
  \item \textbf{Verify} on any SQLite database whether its foreign keys are
        actually being enforced.
\end{tight}
\end{outcomes}

\begin{prereq}
\chref{relational}: relations, keys, nulls. \chref{cardinality} for
participation, which is the constraint this chapter implements. You will need
\texttt{sqlite3} and \texttt{code/m08\_constraints.py}.
\end{prereq}

\begin{keyterms}
integrity constraint $\bullet$ domain constraint $\bullet$ entity integrity
$\bullet$ referential integrity $\bullet$ foreign key $\bullet$ referencing
relation $\bullet$ referenced relation $\bullet$ orphan row $\bullet$
referential action $\bullet$ \sqlkw{CASCADE} $\bullet$ \sqlkw{RESTRICT}
$\bullet$ \sqlkw{SET NULL} $\bullet$ semantic constraint $\bullet$
\sqlkw{CHECK} $\bullet$ \sqlkw{PRAGMA} $\bullet$ declared versus enforced
\end{keyterms}

\section{Three Constraints the Model Gives You Free}

A schema is not only a list of attributes. It is a list of \emph{promises}, and
the front matter of this handout claims that a database is nothing but a set
of promises about data. This is the chapter where that is made literal.

\begin{definitionbox}[The three model-based constraints]
A \term{domain constraint} says every value of an attribute comes from its
domain. \texttt{semester} is an integer between 1 and 8; \texttt{'spring'} is
not a semester and neither is $-3$.

\smallskip
\term{Entity integrity} says no attribute of a primary key may be null. An
identifier that is missing identifies nothing, so the tuple could not be
addressed.

\smallskip
\term{Referential integrity} says that if a tuple refers to another relation,
the tuple it refers to must exist. A section claiming to belong to course
CS-999 when there is no CS-999 is an \term{orphan}.

\smallskip
These three come from the model itself. A fourth kind, \term{semantic}
constraints, come from the university --- ``a student may not enrol in more
than six sections'' --- and have to be stated case by case.
\end{definitionbox}

\begin{examplebox}[One row, three ways to be wrong]
Into \relation{section}{\relkey{course\_code}, \relkey{section\_no}, room,
emp\_id}:

\medskip
\begin{tabular}{@{}L{4.4cm}L{3.2cm}L{5.8cm}@{}}
\textbf{The row} & \textbf{Violates} & \textbf{Why it matters} \\
\texttt{('CS-103','two','R7')} & domain & \texttt{section\_no} is an
  integer; sorting and counting now fail \\
\texttt{(NULL,1,'R7')} & entity integrity & part of the primary key is
  missing, so the row cannot be addressed \\
\texttt{('CS-999',1,'R7')} & referential integrity & there is no course
  CS-999; the timetable now lists a section of nothing \\
\end{tabular}

\smallskip
All three are refusals a database can make \emph{for you}, on every insert,
for ever --- which is exactly what the application code in
\chref{whynotfiles} could not do, because each program had to remember to
check and one of them always forgot.
\end{examplebox}

\section{Foreign Keys}

\begin{definitionbox}[Foreign key]
A \term{foreign key} is a set of attributes $FK$ in a \term{referencing}
relation $R$ whose values must either
\begin{tight}
  \item appear as the primary key of some tuple in the \term{referenced}
        relation $S$, or
  \item be entirely null.
\end{tight}

\smallskip
The null case is not a loophole --- it is \chref{cardinality}'s partial
participation. A section with no instructor yet has a null \texttt{emp\_id},
and that is a design decision you already took.
\end{definitionbox}

\dsfig{%
\begin{tikzpicture}[
  hd/.style={font=\scriptsize\bfseries, text=white},
  cell/.style={font=\scriptsize\ttfamily, text=primary!65!black},
  bad/.style={font=\scriptsize\ttfamily, text=accent},
  ttl/.style={font=\scriptsize\bfseries, text=primary},
  note/.style={font=\scriptsize\itshape, text=neutral, align=left},
  ok/.style={-{Stealth[length=2mm]}, greenhl, line width=0.9pt},
  no/.style={-{Stealth[length=2mm]}, accent, line width=0.9pt, dashed}
]
  % --- parent: course, top right so the arrows have clear air on the left
  \node[ttl, anchor=west] at (2.2,1.42) {course \ \ (referenced)};
  \fill[primary, rounded corners=2pt] (2.2,0.55) rectangle (6.6,1.1);
  \node[hd, anchor=west] at (2.32,0.83) {course\_code};
  \node[hd, anchor=west] at (4.75,0.83) {title};
  \node[cell, anchor=west] at (2.32,0.3)  {CS-103};
  \node[cell, anchor=west] at (4.75,0.3)  {Databases};
  \node[cell, anchor=west] at (2.32,-0.1) {CS-201};
  \node[cell, anchor=west] at (4.75,-0.1) {Data Struct};
  \draw[primary!35, line width=0.8pt] (2.2,-0.33) rectangle (6.6,1.1);
  \node[note, anchor=north west, text=primary] at (2.2,-0.40)
       {course\_code is the primary key};

  % --- child: section
  \node[ttl, anchor=west] at (0,-1.18) {section \ \ (referencing)};
  \fill[accent!10] (0,-3.3) rectangle (6.6,-2.92);
  \fill[primary, rounded corners=2pt] (0,-2.05) rectangle (6.6,-1.5);
  \node[hd, anchor=west] at (0.12,-1.78) {course\_code};
  \node[hd, anchor=west] at (2.6,-1.78)  {section\_no};
  \node[hd, anchor=west] at (4.8,-1.78)  {emp\_id};
  \node[cell, anchor=west] at (0.12,-2.3) {CS-103};
  \node[cell, anchor=west] at (2.6,-2.3)  {1};
  \node[cell, anchor=west] at (4.8,-2.3)  {E-114};
  \node[cell, anchor=west] at (0.12,-2.7) {CS-103};
  \node[cell, anchor=west] at (2.6,-2.7)  {2};
  \node[cell, anchor=west, text=neutral] at (4.8,-2.7) {NULL};
  \node[bad, anchor=west] at (0.12,-3.11) {CS-999};
  \node[bad, anchor=west] at (2.6,-3.11)  {1};
  \node[bad, anchor=west] at (4.8,-3.11)  {E-114};
  \draw[primary!35, line width=0.8pt] (0,-3.33) rectangle (6.6,-1.5);

  % the two valid references, up the left-hand side
  \draw[ok] (0.08,-2.3) .. controls (-1.5,-1.9) and (-1.5,0.5) .. (2.12,0.3);
  \draw[ok] (0.08,-2.7) .. controls (-1.9,-2.1) and (-1.9,0.7) .. (2.12,0.24);
  % The orphan gets a mark beside its row, not an arrow. A dashed arrow
  % "pointing at nothing" had to curve round the outside of the figure and
  % crossed both annotations on the way.
  \node[bad, anchor=west] at (6.75,-3.11) {$\times$ no such course};

  % Both notes below the tables, side by side: nothing can overlap a row.
  \node[note, anchor=north west, text width=5.0cm] at (0,-3.62)
       {\texttt{NULL} is allowed by the definition --- it is Chapter~6's
        partial participation, not a hole in the rule};
  \node[note, anchor=north west, text=accent, text width=5.4cm] at (5.4,-3.62)
       {\textbf{the orphan.} CS-999 does not exist.
        \textbf{Measured: SQLite accepts this row by default.}};
  \draw[-{Stealth[length=1.8mm]}, neutral!60] (1.9,-3.58) -- (4.75,-2.7);
\end{tikzpicture}}%
{Referential integrity. Two of the three child rows point at a course that
exists; the null is permitted by the definition; the third is an orphan. The
whole of this chapter's measurement is that the third row goes in.}{fkorphan}

\begin{conceptbox}[Where foreign keys come from]
You do not invent them. Each one is a \chref{cardinality} constraint written
down in a different notation, and \chref{mapping} turns one into the other
mechanically.

\begin{tight}
  \item \textbf{A 1:N relationship} puts a foreign key on the \emph{many} side,
        pointing at the one. \ent{course} $\rightarrow$
        \ent{department}.
  \item \textbf{Total participation} makes that foreign key \sqlkw{NOT NULL}.
        Partial participation leaves it nullable.
  \item \textbf{A weak entity} gets a foreign key to its owner, and that
        foreign key is part of its own primary key.
  \item \textbf{An M:N relationship} has no foreign key, because it has no
        table yet. It becomes one in \chref{mapping}, with two.
\end{tight}

\smallskip
So the four answers you argued about in \chref{cardinality} are not an
academic exercise: each one lands as a specific line of \sqlkw{CREATE TABLE},
and you can check your diagram against your schema by reading them back.
\end{conceptbox}

\begin{definitionbox}[Referential actions]
Referential integrity can be broken from the \emph{other} end: delete the
course and its sections become orphans. A \term{referential action} says what
the database should do about it.

\medskip
\begin{tabular}{@{}L{3.1cm}L{5.0cm}L{5.9cm}@{}}
\textbf{Action} & \textbf{On deleting the parent} & \textbf{Use when} \\
\sqlkw{RESTRICT} / \sqlkw{NO ACTION} & refuse the delete & the default, and
  the right answer when you are unsure \\
\sqlkw{CASCADE} & delete the children too & the child is meaningless alone ---
  a weak entity \\
\sqlkw{SET NULL} & null the children's foreign key & the child survives
  without the parent, and the column is nullable \\
\end{tabular}

\smallskip
The same three apply \sqlkw{ON UPDATE}, for when a primary key value changes
--- which is one more reason \chref{ermodel} said to choose a key that never
changes.
\end{definitionbox}

\begin{alertbox}[\sqlkw{CASCADE} is a loaded weapon]
\sqlkw{ON DELETE CASCADE} is attractive because it makes a tidy-up problem
disappear. It also means that one \sqlkw{DELETE} statement can remove rows
from tables you did not name, in tables you did not think about, through
chains you did not draw.

\smallskip
Delete a department with cascade all the way down and you remove its courses,
their sections, every enrolment in them, and therefore the grade history of
students who graduated years ago. One statement, no warning, and the
transaction commits.

\smallskip
\textbf{The rule used in this course:} cascade only to a \emph{weak} entity ---
one that genuinely cannot exist without its owner --- and use \sqlkw{RESTRICT}
everywhere else. A refused delete is an inconvenience; a successful one is
sometimes a disaster, and only one of the two is reversible.
\end{alertbox}

\section{Declared Is Not Enforced}

Everything above is standard and true of every relational system. What follows
is specific to the tool on your machine, and it is the most important single
fact in this chapter.

\begin{measured}[A correct \sqlkw{REFERENCES} clause that does nothing]
\textsf{\footnotesize \texttt{code/m08\_constraints.py}, SQLite 3.51.2,
Intel i5-6600, 16\,GB, Windows 11.}

\medskip
\textbf{1. The same schema, the same \sqlkw{INSERT}, two connections}

\begin{tabular}{@{}L{5.6cm}R{2.4cm}L{5.6cm}@{}}
\textbf{\sqlkw{PRAGMA} foreign\_keys} & \textbf{orphans in} & \textbf{outcome} \\
default (untouched) & 1 & \textbf{accepted} \\
\sqlkw{ON}          & 0 & refused: \texttt{FOREIGN KEY constraint failed} \\
\end{tabular}

\smallskip
The \sqlkw{REFERENCES} clause is identical in both. Only the connection
differs.

\medskip
\textbf{2. The setting is per connection, not per database}

\begin{tabular}{@{}L{8.6cm}L{5.0cm}@{}}
connection A, after \sqlkw{PRAGMA} foreign\_keys = \sqlkw{ON} & reports 1 \\
connection A, orphan insert & refused \\
connection B, \emph{same file}, nothing set & reports 0 \\
connection B, orphan insert & \textbf{accepted} \\
\end{tabular}

\smallskip
Nothing in the file remembers that enforcement was ever wanted.

\medskip
\textbf{3. A \sqlkw{PRIMARY KEY} column that stores a null}

\begin{tabular}{@{}L{5.6cm}R{2.6cm}R{2.4cm}R{2.0cm}@{}}
\textbf{declared as} & \textbf{\sqlkw{INSERT} NULL} & \textbf{stored} & \textbf{rows} \\
\texttt{TEXT PRIMARY KEY}            & accepted & \textbf{null} & 1 \\
\texttt{INTEGER PRIMARY KEY}         & accepted & integer & 1 \\
\texttt{TEXT PRIMARY KEY}, \sqlkw{STRICT} & refused & --- & 0 \\
\texttt{TEXT PRIMARY KEY NOT NULL}   & refused & --- & 0 \\
\end{tabular}

\smallskip
Read the \emph{stored} column, not the first one. \texttt{INTEGER PRIMARY KEY}
accepted the null and stored an integer --- it is an alias for the rowid, so
the null was replaced, and no constraint was broken. \texttt{TEXT PRIMARY KEY}
accepted it and \textbf{stored null}, which breaks entity integrity --- and
that is the declaration every table in this course uses.

\medskip
\textbf{4. What enforcement costs} (20{,}000 valid child rows, 5 runs)

\begin{tabular}{@{}L{6.6cm}R{2.4cm}R{2.2cm}R{2.2cm}@{}}
& \textbf{median} & \textbf{min} & \textbf{max} \\
foreign\_keys \sqlkw{OFF} & 58.7\,ms & 55.1 & 63.8 \\
foreign\_keys \sqlkw{ON}  & 73.0\,ms & 67.3 & 79.6 \\
\textbf{cost of enforcement} & \textbf{14.3\,ms} & \multicolumn{2}{l}{\textbf{(+24\%)}} \\
\end{tabular}

\smallskip
\textbf{Read the four results together.} The declaration is accepted, stored,
and printed back by \texttt{.schema} --- and by default enforces nothing. It
costs 24\% on bulk inserts to turn on, which is a real cost and a small one:
note that the slowest run with enforcement on (79.6\,ms) is only 25\% worse
than the slowest without it, and that the two ranges do not overlap, so the
24\% is a result rather than noise. \textbf{A single run of this measurement
put the two settings 0.6\,ms apart}, which would have been reported as 3\%.
Repeating it is not a nicety.
\end{measured}

\begin{alertbox}[What this means for everything you write from here on]
\sqlkw{PRAGMA} \texttt{foreign\_keys} has been off by default since foreign
keys were added to SQLite in 2009, and it will stay off: turning it on would
break existing databases that have been quietly accumulating orphans.

\smallskip
The consequences you have to live with:

\begin{tight}
  \item \textbf{Set it on every connection}, as the first statement, before
        any transaction. It is a no-op inside one.
  \item \textbf{Your lab partner's connection is not yours.} A database can be
        enforced in your program and unenforced in theirs, at the same moment,
        on the same file.
  \item \textbf{Orphans already in the file are not detected} when you turn it
        on. Enforcement applies to \emph{new} statements only.
        \texttt{PRAGMA foreign\_key\_check} is how you find the old ones, and
        step 7 of the laboratory runs it.
  \item \textbf{A schema that looks correct may never have been enforced.}
        When you are handed an unfamiliar SQLite database, this is the first
        thing to check and the last thing anybody remembers.
\end{tight}

\smallskip
This is not a criticism of SQLite --- the behaviour is documented and the
reason is sound. It is a demonstration of something general: \textbf{a
constraint has two separate lives, one in the schema and one at run time, and
only the second one protects anything.}
\end{alertbox}

\begin{examplebox}[The university schema, with its promises declared]
\begin{lstlisting}[language=sql]
PRAGMA foreign_keys = ON;        -- first. every time.

CREATE TABLE department (
  dept_code TEXT PRIMARY KEY NOT NULL,
  name      TEXT NOT NULL
) STRICT;
\end{lstlisting}

A course belongs to exactly one department, and must:

\begin{lstlisting}[language=sql]
CREATE TABLE course (
  course_code  TEXT PRIMARY KEY NOT NULL,
  title        TEXT NOT NULL,
  credit_hours INTEGER NOT NULL CHECK (credit_hours BETWEEN 1 AND 4),
  dept_code    TEXT NOT NULL              -- total participation
               REFERENCES department(dept_code)
               ON DELETE RESTRICT         -- a department with courses stays
) STRICT;
\end{lstlisting}

A section is the weak entity of \chref{cardinality}, so its owner's key is part
of its own and the cascade is correct here and nowhere else:

\begin{lstlisting}[language=sql]
CREATE TABLE section (
  course_code TEXT    NOT NULL
              REFERENCES course(course_code)
              ON DELETE CASCADE,          -- weak entity: cascade is correct
  semester    TEXT    NOT NULL,
  section_no  INTEGER NOT NULL CHECK (section_no > 0),
  room        TEXT,
  emp_id      TEXT                        -- nullable: partial participation
              REFERENCES instructor(emp_id)
              ON DELETE SET NULL,         -- the section survives; the
  PRIMARY KEY (course_code, semester, section_no)   -- teacher changed
) STRICT;
\end{lstlisting}

\smallskip
\textbf{Every line is a decision you have already made.} \texttt{NOT NULL} on
\texttt{dept\_code} is the total participation from \chref{cardinality}'s
worked example; the nullable \texttt{emp\_id} is the partial participation
argued for in the same place; the three-part primary key is the weak entity's
owner key plus discriminator; and each of the three referential actions
follows the rule in this chapter. Nothing here is a matter of taste.
\end{examplebox}

\begin{pitfall}
\begin{tight}
  \item \textbf{Assuming \sqlkw{REFERENCES} does something.} By default, in
        SQLite, it does not. Check, every time.
  \item \textbf{Setting the pragma once and assuming it stuck.} It is per
        connection and lasts until that connection closes.
  \item \textbf{Setting it inside a transaction.} It is silently ignored there.
  \item \textbf{Declaring \texttt{TEXT PRIMARY KEY} without \sqlkw{NOT NULL}.}
        It will store a null, as measured. Add \sqlkw{NOT NULL} or
        \sqlkw{STRICT}, and this course does both.
  \item \textbf{Cascading beyond a weak entity.} One delete, a chain you did
        not draw, and a committed transaction.
  \item \textbf{Believing the application will check it.} Three programs, one
        of them written next year by somebody else. This is
        \chref{whynotfiles}'s whole argument.
  \item \textbf{Turning enforcement on and assuming the old data is clean.} It
        applies to new statements only; run
        \texttt{PRAGMA foreign\_key\_check}.
\end{tight}
\end{pitfall}

\begin{lab}[Insert an orphan into a correct schema, then make it impossible]
\begin{lstlisting}[language=bash]
# --- 1. A scratch database ------------------------------------------
sqlite3 fk.db
\end{lstlisting}
\begin{lstlisting}[language=sql]
-- --- 2. A schema that declares referential integrity ------------
.headers on
.mode box
CREATE TABLE course (
  course_code TEXT PRIMARY KEY NOT NULL,
  title       TEXT NOT NULL
) STRICT;
CREATE TABLE section (
  course_code TEXT    NOT NULL REFERENCES course(course_code),
  section_no  INTEGER NOT NULL,
  room        TEXT,
  PRIMARY KEY (course_code, section_no)
) STRICT;
INSERT INTO course VALUES ('CS-103','Database Systems');

-- --- 3. Confirm the declaration is really there -----------------
.schema section
PRAGMA foreign_key_list(section);
-- It is declared. Written down. Reported back to you.

-- --- 4. Now break it --------------------------------------------
PRAGMA foreign_keys;              -- what does it say?
INSERT INTO section VALUES ('CS-999', 1, 'Room 7');
SELECT * FROM section;
-- A section of a course that does not exist. No error.

-- --- 5. Turn enforcement on and try again -----------------------
PRAGMA foreign_keys = ON;
PRAGMA foreign_keys;
INSERT INTO section VALUES ('CS-888', 1, 'Room 8');
-- Record the exact error message.

-- --- 6. The orphan from step 4 is still there -------------------
SELECT * FROM section;
-- Enforcement did not clean up. It never looks backwards.

-- --- 7. Find every orphan already in the file -------------------
PRAGMA foreign_key_check;
-- Learn this one. It is how you audit a database you were handed.

-- --- 8. Prove it is per connection ------------------------------
.quit
\end{lstlisting}
\begin{lstlisting}[language=bash]
# --- 9. A brand new connection to the SAME file ---------------------
sqlite3 fk.db "PRAGMA foreign_keys;"
sqlite3 fk.db "INSERT INTO section VALUES ('CS-777',1,'R9');"
sqlite3 fk.db "SELECT * FROM section;"
# It went in. Nothing in the file remembered.

# --- 10. The measured version ---------------------------------------
python code/m08_constraints.py
\end{lstlisting}

\smallskip
\textbf{What to notice.} Step 3 is the one that makes the point. The constraint
is \emph{there} --- \texttt{PRAGMA foreign\_key\_list} prints it, the schema
stores it, any tool reading the database would report it --- and step 4 walks
straight past it. From now on, every script you write in this course opens
with \texttt{PRAGMA foreign\_keys = ON;}, and \texttt{code/run\_all.py}
checks that they all do.
\end{lab}

\begin{checkpoint}
\begin{tightnum}
  \item You are handed a SQLite database and told its foreign keys are
        enforced. List the commands you would run to find out whether that is
        true \emph{now}, and whether it has ever been true, and say what each
        result would and would not establish.
  \item For each of \ent{section}~$\rightarrow$~\ent{course},
        \ent{section}~$\rightarrow$~\ent{instructor} and
        \ent{enrolment}~$\rightarrow$~\ent{student}, choose an
        \sqlkw{ON DELETE} action and justify it in one sentence. At least one
        of your three must not be \sqlkw{CASCADE}.
  \item A table is declared \texttt{roll\_no TEXT PRIMARY KEY}. Explain what
        value it will accept that it should not, which integrity constraint
        that violates, and give the two different one-word changes that each
        fix it.
\end{tightnum}
\end{checkpoint}

\begin{chaptersummary}
\begin{tight}
  \item The model gives three constraints: \textbf{domain} (values come from
        the right set), \textbf{entity integrity} (no part of a primary key is
        null), \textbf{referential integrity} (a reference points at something
        that exists). \textbf{Semantic} constraints come from the university.
  \item A \textbf{foreign key} must match a primary key in the referenced
        relation, \emph{or} be wholly null --- and the null case is
        \chref{cardinality}'s partial participation.
  \item Foreign keys are not invented: a 1:N relationship puts one on the many
        side, total participation makes it \sqlkw{NOT NULL}, and a weak entity
        makes it part of its own key.
  \item \textbf{Referential actions}: \sqlkw{RESTRICT} by default,
        \sqlkw{CASCADE} only to a weak entity, \sqlkw{SET NULL} when the child
        survives the parent.
  \item \textbf{Measured: \sqlkw{PRAGMA} foreign\_keys is 0 by default.} A
        correct \sqlkw{REFERENCES} clause is parsed, stored, reported --- and
        enforces nothing. The orphan went in without an error.
  \item The setting is \textbf{per connection}. Connection B, on the same
        file, accepted what connection A refused.
  \item Enforcement \textbf{never looks backwards}: orphans already present
        stay. \texttt{PRAGMA foreign\_key\_check} finds them.
  \item \textbf{Measured: \texttt{TEXT PRIMARY KEY} stores a null.}
        \sqlkw{STRICT} or \sqlkw{NOT NULL} fixes it; \texttt{INTEGER PRIMARY
        KEY} only appears to, because it substitutes a rowid.
  \item \textbf{Measured: enforcement costs +24\%} on bulk inserts --- outside
        the run-to-run spread, and cheap for what it buys. A single run
        reported 3\%, which is why the harness repeats.
  \item The general lesson: \textbf{a constraint has two lives, one in the
        schema and one at run time, and only the second protects anything.}
\end{tight}
\end{chaptersummary}

\begin{reviewq}
  \item \textbf{(MCQ)} Entity integrity is violated when: (a)~a foreign key
        points nowhere; (b)~part of a primary key is null; (c)~a value is
        outside its domain; (d)~two tables have the same name.
  \item \textbf{(MCQ)} In SQLite, \sqlkw{PRAGMA} \texttt{foreign\_keys}
        defaults to: (a)~\sqlkw{ON} for \sqlkw{STRICT} tables;
        (b)~\sqlkw{ON}; (c)~\sqlkw{OFF}; (d)~whatever it was last set to in
        that file.
  \item \textbf{(Short)} Give a situation in which \sqlkw{ON DELETE SET NULL}
        is the right action and \sqlkw{CASCADE} would lose data you need.
  \item \textbf{(Short)} Why is a foreign key allowed to be null at all? Answer
        in terms of \chref{cardinality}.
  \item \textbf{(Short)} You turn on foreign key enforcement for a database
        that has been running for a year. What does it protect you from, and
        what does it not?
  \item \textbf{(Applied)} Write the \sqlkw{CREATE TABLE} statements for
        \ent{instructor} and \ent{enrolment}, with every domain, entity
        and referential constraint declared and a referential action chosen
        for each foreign key. Beside each constraint, name the decision from
        \chref{ermodel} or \chref{cardinality} that it implements.
  \item \textbf{(Applied)} A colleague argues that foreign keys should be left
        off because they cost 24\% on bulk inserts and ``the application
        checks anyway''. Write the reply. Use the measured figures, the
        per-connection result, and \chref{whynotfiles}'s argument, and say
        under what single condition their position would actually be
        defensible.
\end{reviewq}
